\chapter{Commodities and Energy}\label{ch:iv2:commodities-and-energy}

\begin{ivtrains}
Curve shapes and what they say about storage, spreads between products and between places, seasonal curves and power, and the options that a physical asset really is. Every level here is invented for the arithmetic, except where a box gives a date.
\end{ivtrains}

\begin{question}[$\star$ \iqroles{Trader, QT, Analyst} \iqfirm{any} \iqlevel{junior}]\label{iq:iv2:commodities-and-energy:1}
Crude oil is 70 in the spot market, and the one-year future is 74. What is this curve shape called? You buy the future. If the spot price is still 70 in a year, what do you make or lose per barrel?
\end{question}

\begin{question}[$\star$ \iqroles{Trader, QT, QR} \iqfirm{any} \iqlevel{junior}]\label{iq:iv2:commodities-and-energy:2}
A commodity costs 80 today, the rate is 5\%, and storage costs 4 for the year, paid at the end. What is the highest the one-year future can be? It trades at 92, and you have access to storage. What do you do, and how much do you lock in?
\end{question}

\begin{question}[$\star$ \iqroles{Trader, Analyst, QT} \iqfirm{hedge fund, bank} \iqlevel{junior}]\label{iq:iv2:commodities-and-energy:3}
Crude is \$80 a barrel, gasoline \$2.40 a gallon, heating oil \$2.70 a gallon, and a barrel holds 42 gallons. What is the 3:2:1 crack spread per barrel of crude? Who is naturally long it?
\end{question}

\begin{question}[$\star$ \iqroles{Trader, QT} \iqfirm{hedge fund, bank} \iqlevel{junior}]\label{iq:iv2:commodities-and-energy:4}
Front-month crude is 72, and the second month is 70. Each contract is 1,000 barrels. You buy the front and sell the second. What do you make if both rise by \$5? If the front rises \$3 and the second \$2? What view are you expressing?
\end{question}

\begin{question}[$\star$ \iqroles{Trader, Analyst, QT} \iqfirm{bank, hedge fund} \iqlevel{junior}]\label{iq:iv2:commodities-and-energy:5}
Power trades at \$60 per MWh and gas at \$5 per MMBtu, and a plant has a heat rate of 7 MMBtu per MWh. What is the spark spread? How much does a 100 MW plant earn in a 16-hour block? With carbon at \$30 a tonne and 0.4 tonnes per MWh, what is the clean spark spread?
\end{question}

\begin{question}[$\star\star$ \iqroles{QR, Trader, QT} \iqfirm{any} \iqlevel{junior, experienced}]\label{iq:iv2:commodities-and-energy:6}
A metal has spot 70 and a one-year future at 66. The rate is 4\% and storage costs 2\% of value a year, both continuous. What convenience yield does that imply? What does the number mean, and when would it be high?
\end{question}

\begin{question}[$\star\star$ \iqroles{Trader, QR, QT} \iqfirm{any} \iqlevel{junior, experienced}]\label{iq:iv2:commodities-and-energy:7}
Spot is 50, the rate is 4\% and storage is \$1 a month. How wide can the one-month contango be? Is there a matching limit on how far backwardation can go? Why or why not?
\end{question}

\begin{question}[$\star\star$ \iqroles{QR, QT, Trader} \iqfirm{hedge fund, bank} \iqlevel{junior, experienced}]\label{iq:iv2:commodities-and-energy:8}
Commodity A's curve slopes up by 1\% a month, forever, and its spot price never moves. Gas has a purely seasonal curve: each calendar month has a price of $3 + 0.6\cos(2\pi(m - 1)/12)$, and each contract expires exactly at that price. What does a long position that rolls monthly earn in a year in each? How much of the gas curve's slope is a real return?
\end{question}

\begin{question}[$\star\star$ \iqroles{Trader, QT, QR} \iqfirm{hedge fund, bank} \iqlevel{junior}]\label{iq:iv2:commodities-and-energy:9}
An airline will buy 10 million gallons of jet fuel in three months and hedges with heating oil futures of 42,000 gallons each. The monthly volatility of jet fuel price changes is 3\%, that of heating oil is 3.5\%, and their correlation is 0.8. What hedge ratio and how many contracts? How much of the variance does it remove?
\end{question}

\begin{question}[$\star\star$ \iqroles{QR, Trader} \iqfirm{hedge fund, bank} \iqlevel{experienced}]\label{iq:iv2:commodities-and-energy:10}
A future with time to maturity $\tau$ has volatility $0.6\,e^{-2\tau}$. What is the volatility of the one-month future today, and of the one-year future? What average volatility would you use for a one-year option on the one-year future? For a one-month option on the one-month future? What is the economic story?
\end{question}

\begin{question}[$\star\star$ \iqroles{QT, QR, Trader} \iqfirm{hedge fund, bank} \iqlevel{junior, experienced}]\label{iq:iv2:commodities-and-energy:11}
A one-month at-the-money option on a future at 70, with 40\% volatility, is settled against the average of the price over the whole month. About what is it worth, and what is a plain European option worth? Why do so many commodity options settle this way?
\end{question}

\begin{question}[$\star\star$ \iqroles{Trader, Analyst, QT} \iqfirm{hedge fund, bank} \iqlevel{junior, experienced}]\label{iq:iv2:commodities-and-energy:12}
US crude costs 73 at the port. A cargo takes 30 days to reach Europe, and freight costs \$3.50 a barrel. The European benchmark for delivery in 30 days is 77. Financing is 5\%. Is the arbitrage open? How do you lock it in, and what risk remains?
\end{question}

\begin{question}[$\star\star$ \iqroles{Trader, QT, Analyst} \iqfirm{bank, hedge fund} \iqlevel{junior, experienced}]\label{iq:iv2:commodities-and-energy:13}
A power market has four sets of plants: wind 30 MW at \$0, nuclear 40 MW at \$10, combined-cycle gas 40 MW at \$40 and peakers 20 MW at \$90. What is the price at demand of 20, 60, 100 and 115 MW? What happens at 140? Wind capacity doubles to 60 MW. What happens to the prices? Why can power prices go negative?
\end{question}

\begin{question}[$\star\star\star$ \iqroles{QR, Trader} \iqfirm{bank, hedge fund} \iqlevel{experienced}]\label{iq:iv2:commodities-and-energy:14}
The power forward is 40 per MWh and the gas forward is 5 per MMBtu. A plant has a heat rate of 7 MMBtu per MWh. Power volatility is 60\%, gas volatility 40\%, their correlation 0.5, and the option to run expires in three months. Ignoring discounting, value the option to run for one MWh with Margrabe. What is the intrinsic value, and what is the rest? What happens at a correlation of 0.9? How do you hedge it?
\end{question}

\begin{question}[$\star\star\star$ \iqroles{Trader, QR, Risk} \iqfirm{hedge fund, bank} \iqlevel{junior, experienced}]\label{iq:iv2:commodities-and-energy:15}
A crude future for physical delivery has one day left. The next month trades at 20. The storage at the delivery point is effectively full. What stops the expiring contract from trading below zero? At what storage cost, per barrel per month, does a price of $-38$ make sense? Who has to sell at any price?
\end{question}

\begin{question}[$\star\star\star$ \iqroles{QR, Trader, QD} \iqfirm{bank, hedge fund} \iqlevel{junior, experienced}]\label{iq:iv2:commodities-and-energy:16}
A gas storage tank holds one unit and can inject or withdraw one unit a month. The forward curve for five months is 3.00, 2.90, 2.85, 3.40, 3.60, and the tank ends empty. Injection costs 0.02 and withdrawal 0.01. What is the intrinsic value? Then month 5's price is revealed in month 4, at 4.10 or 3.10 with equal chance. What is the tank worth when you may withdraw in month 4 instead?
\end{question}

\begin{question}[$\star\star\star$ \iqroles{Trader, Risk, QR} \iqfirm{bank, hedge fund} \iqlevel{experienced}]\label{iq:iv2:commodities-and-energy:17}
You have sold 10 million barrels at a fixed price, 1 million barrels a year for ten years, and hedge one for one with long front-month futures, rolled each month. Oil falls \$20 in a flat move across the curve, and the rate is 5\%. What are your margin call and the present value of your gain on the sales? Is this hedge right? Now the curve moves only 30\% of the way at the long end, or goes into a contango of 50 cents a month. What then?
\end{question}

\begin{question}[$\star\star\star$ \iqroles{Trader, QR, Analyst} \iqfirm{bank, hedge fund} \iqlevel{junior, experienced}]\label{iq:iv2:commodities-and-energy:18}
A producer wants to protect a one-year future at 70, with no premium, using a put at 60 and a call. Rates are zero and volatility is flat at 30\%. What call strike makes it costless? With skew, the put is at 36\% volatility and the call at 26\%. What now? What does the producer give up?
\end{question}

\begin{question}[$\star\star\star$ \iqroles{QR, Trader} \iqfirm{bank, hedge fund} \iqlevel{experienced}]\label{iq:iv2:commodities-and-energy:19}
You price a six-month option on a refining spread: payoff $\max(F_1 - F_2 - 15, 0)$ with $F_1 = 100$, $F_2 = 80$, volatilities 35\% and 30\%, and correlation 0.9. Use Kirk's approximation. Then give a second estimate treating the spread as normal. Which would you trust, and how would you check it?
\end{question}

\begin{question}[$\star\star\star\star$ \iqroles{QR, Trader} \iqfirm{bank, hedge fund} \iqlevel{experienced}]\label{iq:iv2:commodities-and-energy:20}
You are offered a tolling deal for one month: you pay a fee, and can run a 100 MW plant, heat rate 8, on 22 days for 16 hours a day, buying gas and selling power. The power forward is \$42, the gas forward \$5. Monthly forward volatilities are 50\% and 35\%, correlation 0.6, and the month is one month away. Day-ahead volatilities are 90\% and 50\%, correlation 0.5, and take each day's option as expiring in $1/8$ of a year. What is the intrinsic value? What is it worth as one monthly decision, and as 22 daily ones? What hedge do you put on against each? What fee would you pay, and why not the top of your range? What happens to the value if the forward is \$32?
\end{question}
